(1) 取回路电荷 \(q\) 和电流 \(i=\dot q\)，基尔霍夫电压定律给 \(L\ddot q+R\dot q+q/C(t)=V(t)\)。电容端电压瞬时为 \(q/C(t)\)，但因为电容随时间变化，电容电流不能错误写成 \(C(t)\dot V_C\) 而遗漏 \(\dot C V_C\)。此式同时覆盖外部电压激励与纯参数调制两种功率通道。
(2) 总储能 \(U=Li^2/2+q^2/[2C(t)]\)，求导并代入回路方程，得到 \(\dot U=Vi-Ri^2-q^2\dot C/(2C^2)\)。末项是外部调整电容几何或介质时向电路注入的功率，符号取决于调制阶段；其平均值可以为正。若省略该项，会错误断言无电源回路不可能出现参数放大。
(3) 对 \(h\ll1\)，\(1/C(t)=C_0^{-1}[1-h\cos(2\omega_0t)+O(h^2)]\)，无源方程近似为 \(\ddot q+2\gamma\dot q+\omega_0^2[1-h\cos(2\omega_0t)]q=0\)，\(\gamma=R/(2L)\)。令 \(q=A(t)\cos\omega_0t+B(t)\sin\omega_0t\) 并在一个快周期内平均，慢幅度的最大实增长率为 \(\sigma=h\omega_0/4-\gamma+O(h^2,\gamma^2/\omega_0)\)，对应二倍频参数泵浦。
(4) 增益要求 \(\sigma>0\)，所以一阶阈值为 \(h>4\gamma/\omega_0=2R/(L\omega_0)\)；若 \(R=0\)，理想无耗散系统在精确共振处任意非零小调制都可激发增长。真实回路存在电容调节范围、放大后非线性、介质损耗与泵浦功率有限等限制，指数律仅是线性小信号稳定性结果。
